\documentclass{article}

% "preprint" option is used for arXiv or other preprint submissions
\usepackage[preprint]{neurips_2025}

\usepackage[utf8]{inputenc}
\usepackage[T1]{fontenc}
\usepackage[spanish]{babel}
\usepackage{hyperref}
\usepackage{url}
\usepackage{booktabs}
\usepackage{amsfonts}
\usepackage{amsmath}
\usepackage{nicefrac}
\usepackage{microtype}
\usepackage{xcolor}
\usepackage{array}
\usepackage{multirow}

\title{No Free Lunch: Hedge Fund Time Series Forecasting}

\author{%
  Andrés Mosquera-Hernández \quad Josue Briceño Urquijo \quad Michael Patiño Pantoja \\
  Universidad de los Andes, Bogotá, Colombia \\
}

\begin{document}

\maketitle

% \begin{abstract}
%   TODO: escribir el abstract una vez finalizado el trabajo.
% \end{abstract}


% ============================================================
% 1. INTRODUCCIÓN
% ============================================================

\section{Introducción y Justificación}

El pronóstico de series de tiempo es fundamental en las finanzas cuantitativas —base de la asignación de portafolios, gestión de riesgo y trading sistemático—. Las competencias M4 \citep{makridakis2020m4} y M5 \citep{makridakis2022m5} evidenciaron que los árboles de gradient boosting \citep{ke2017lightgbm} y las arquitecturas híbridas superan sistemáticamente las líneas base estadísticas clásicas, y las soluciones ganadoras en Kaggle confirman que la validación cruzada temporal y los ensambles recursivo-directos son decisivos para la generalización fuera de muestra \citep{m5first2020, kaggleguide}. En la práctica, las series financieras añaden dificultades propias: ratio señal-ruido extremadamente bajo, no estacionariedad y dinámicas heterogéneas. La competencia \emph{Hedge Fund Time Series Forecasting} \citep{adatacompany2026} agudiza este escenario al prohibir información externa y penalizar severamente la fuga de datos, siendo el reto principal generalizar sin sobreajustarse a un leaderboard público ruidoso.

Este trabajo se justifica en tres dimensiones: (i)~\textbf{relevancia práctica}, pues mejoras marginales en precisión se traducen en retornos ajustados por riesgo superiores en fondos reales ($>\$4$ billones USD bajo gestión), generando metodología directamente transferible a producción; (ii)~\textbf{vacío de conocimiento}, ya que la mayoría de investigaciones requieren features interpretables (ej. OHLCV, macroeconomía), mientras que aquí evaluamos si el aprendizaje puramente guiado por los datos puede recuperar estructura predictiva sin conocimiento explícito de dominio; y (iii)~\textbf{pertinencia metodológica}, pues el alto riesgo de sobreajuste al leaderboard exige un framework de validación cruzada temporal escrupuloso y reproducible.

% -----------------------------------------------------------
\subsection{Objetivos}

\textbf{Objetivo general.} Desarrollar una metodología de machine learning para el pronóstico de series de tiempo financieras anonimizadas que alcance un \emph{weighted RMSE skill score} superior a 0.30 en el leaderboard público de Kaggle \emph{Hedge Fund - Time Series Forecasting}.

Para cumplir con el propósito principal, se establecen los siguientes \textbf{objetivos específicos}:

\begin{itemize}
    \item Implementar un pipeline que garantice la restricción $\texttt{ts\_index} \leq t$ en todos los estadísticos, rezagos e imputaciones, con cero instancias de fuga de datos verificadas computacionalmente sobre los conjuntos de entrenamiento y test.
    \item Evaluar si incorporar \texttt{code}, \texttt{sub\_code}, \texttt{sub\_category} y \texttt{horizon} como variables de tipo \texttt{category} mejora el \emph{score} local en al menos un 10\% respecto a un modelo sin dichos identificadores, y si la especialización por \texttt{sub\_category} supera a la por horizonte en términos de \emph{score} de validación.
    \item Comparar al menos dos arquitecturas de redes neuronales (MLP, LSTM) y cuatro variantes de LightGBM bajo validación temporal estricta, seleccionando la configuración con mayor \emph{score} en el leaderboard público.
\end{itemize}


% ============================================================
% 2. ANTECEDENTES Y TRABAJOS RELACIONADOS
% ============================================================

\section{Tres trabajos relevantes}

\paragraph{TimeXer: Empowering Transformers for Time Series Forecasting with Exogenous Variables \citep{wang2024timexer}.}
La mayoría de modelos manejan bien la dependencia temporal o las relaciones inter-variable, pero raramente ambas a la vez, y las variables exógenas con irregularidades (datos faltantes, frecuencias distintas) agravan el problema. TimeXer lo resuelve con una representación jerárquica: la serie objetivo se divide en \textit{patches} temporales (dependencia interna) y las variables exógenas se codifican como tokens de variable; un token global aprendible conecta ambas representaciones vía \textit{cross-attention}. Evaluado en datasets meteorológicos, de tráfico y energía, supera a modelos Transformer, RNN y basados en grafos en escenarios univariados y multivariados, confirmando que el tratamiento explícito de exógenas mejora significativamente el forecasting.

\paragraph{Are Self-Attentions Effective for Time Series Forecasting? \citep{kim2024cats}.}
CATS (Cross-Attention-only Time Series transformer) elimina por completo la auto-atención, argumentando que su invariancia de permutación impide capturar dependencias secuenciales. En su lugar, la entrada se codifica como keys y values mientras los queries son parámetros aprendibles dependientes del horizonte, con compartición de parámetros entre horizontes y patching para reducir costo computacional. El resultado es un nuevo estado del arte en MSE y MAE, superando a TimeMixer y PatchTST en horizontes de 96 a 720 pasos, con mejor captura de periodicidad y shocks abruptos.

\paragraph{LightGBM \citep{ke2017lightgbm}.}
LightGBM introduce dos innovaciones: GOSS (\textit{Gradient-based One-Side Sampling}), que retiene instancias de gradiente alto y submuestrea las de gradiente bajo con factor de amplificación, aproximando las ganancias de split con alta fidelidad sobre una fracción de los datos; y EFB (\textit{Exclusive Feature Bundling}), que agrupa features dispersas mutuamente excluyentes mediante coloración greedy de grafos, reduciendo la dimensionalidad sin pérdida de información. Su limitación estructural es tratar cada fila como observación i.i.d.: no hay mecanismo de causalidad interno, por lo que toda la información secuencial debe inyectarse como features externas. En este trabajo esa restricción es precisamente el diferenciador: LightGBM actúa como motor predictivo global mientras que el pipeline construye todos los estadísticos rolling y rezagos bajo la restricción $\texttt{ts\_index} \leq t$, replicando el resultado de M5 donde esta combinación superó a todos los enfoques estadísticos clásicos \citep{makridakis2022m5, m5first2020}.

% ============================================================
% 3. DATOS Y TÉCNICA DE EXTRACCIÓN
% ============================================================

\section{Datos empleados}

Los datos provienen de la competencia Kaggle \emph{Hedge Fund Time Series Forecasting} \citep{adatacompany2026}, aportados por un fondo de cobertura activo. Se distribuyen en dos archivos Parquet: \texttt{train.parquet} (94 columnas, incluye target y pesos) y \texttt{test.parquet} (92 columnas, sin target ni pesos). Las observaciones de test corresponden a un periodo temporal estrictamente posterior al de entrenamiento, aunque las fechas exactas permanecen ocultas.

\begin{table}[h!]
  \caption{Diccionario de datos. El campo \texttt{id} concatena las cinco columnas clave con separador \texttt{\_\_}.}
  \label{tab:datos}
  \centering
  \begin{tabular}{llp{4.8cm}p{3.0cm}}
    \toprule
    Columna & Tipo & Descripción & Implicación para el modelo \\
    \midrule
    \texttt{id}            & str   & Concatenación: \texttt{code}, \texttt{sub\_code}, \texttt{sub\_category}, \texttt{horizon}, \texttt{ts\_index} & Clave única de submisión. \\
    \texttt{code}          & str   & Entidad primaria (anonimizado) & \multirow{3}{3cm}{Jerarquía inter-grupo explotable.} \\
    \texttt{sub\_code}     & str   & Subgrupo dentro de \texttt{code} & \\
    \texttt{sub\_category} & str   & Categoría amplia (anonimizada) & \\
    \texttt{ts\_index}     & int   & Eje temporal entero; sin calendario explícito & Orden causal estricto. \\
    \texttt{horizon}       & int   & Régimen: 1 (corto), 3 (medio), 10 (largo), 25 (extra-largo) & Etiqueta, \emph{no} desplazamiento. \\
    \texttt{weight}        & float & Peso por muestra para la función de pérdida & Solo \texttt{sample\_weight}; \textbf{nunca} feature. \\
    \texttt{y\_target}     & float & Variable objetivo continua (solo en train) & SNR bajo (${\approx}0.063$ baseline). \\
    \texttt{feature\_a} - \texttt{feature\_ch} & float & 86 features numéricas anónimas; admiten nulos & Imputación causal (\emph{forward-fill}). \\
    \bottomrule
  \end{tabular}
\end{table}


% ============================================================
% 4. METODOLOGÍA
% ============================================================

\section{Metodología}
\label{sec:metodologia}

Se desarrollaron ocho experimentos de forma secuencial; los dos primeros exploran arquitecturas de redes neuronales y los seis restantes iteran sobre LightGBM. Todos comparten la restricción causal de que ningún estadístico derivado puede referenciar $\texttt{ts\_index} > t$.

% -----------------------------------------------------------
\subsection{Experimento MLP: Red neuronal densa con embeddings de entidad}

Se trabaja sobre un submuestreo del 25\% de entidades (\texttt{code}, \texttt{sub\_code}), preservando la evolución temporal completa de cada entidad seleccionada. Las cuatro variables categóricas se codifican con \texttt{LabelEncoder} (ajustado sobre train+test) y se representan con capas de embedding de dimensiones $[8, 16, 4, 4]$ para \texttt{code}, \texttt{sub\_code}, \texttt{sub\_category} y \texttt{horizon} respectivamente. Los 86 features numéricos se normalizan con \texttt{StandardScaler} (ajustado solo sobre train) y los valores nulos se imputan con cero. La arquitectura concatena los embeddings, los features numéricos y el $\texttt{ts\_index}$ normalizado, y los pasa por una red MLP con capas $[2048, 1024, 512, 256, 128, 64]$ con \texttt{BatchNorm}, \texttt{ReLU} y \texttt{Dropout}$(0.3)$. Se usa pérdida \texttt{MSELoss}, optimizador \texttt{Adam} (\texttt{lr}=$10^{-3}$, \texttt{weight\_decay}=$10^{-5}$) con clipping de gradientes (\texttt{max\_norm}=1.0) y \texttt{ReduceLROnPlateau}. El corte de validación es el 85\% del \texttt{ts\_index} máximo; el modelo final se reentrena sobre todos los datos con \texttt{CosineAnnealingLR} a 50 épocas y lotes de 4096.

% -----------------------------------------------------------
\subsection{Experimento LSTM: Red recurrente con embeddings de entidad}

Mismo submuestreo del 25\% de entidades y mismo preprocesamiento que el MLP. Los 86 features numéricos se tratan como una secuencia de un solo paso temporal. Cada variable categórica obtiene una capa de embedding de dimensión 86; los cinco tensores (secuencia + 4 embeddings) se concatenan a lo largo del eje de secuencia y se procesan con una capa \texttt{LSTM}(256), seguida de \texttt{Dense}(128, \texttt{relu}) $\to$ \texttt{Dropout}(0.2) $\to$ \texttt{Dense}(64, \texttt{relu}) $\to$ \texttt{Dense}(1). La función de pérdida es \texttt{MAE}, el optimizador \texttt{Adam} (\texttt{lr}=$10^{-5}$) y el scheduler \texttt{ReduceLROnPlateau} (factor=0.5, paciencia=10). El entrenamiento se realiza con hasta 100 épocas, lotes de 256 y datos mezclados en cada época.

% -----------------------------------------------------------
\subsection{Experimento 1: LightGBM multi-horizonte con ingeniería de features causales}

Se entrena un modelo LightGBM independiente por horizonte $h \in \{1,3,10,25\}$, construyendo 183 estadísticos causales organizados en seis grupos:

\begin{enumerate}
  \item \textbf{Identidad}: dummies de \texttt{sub\_category}; codificación de frecuencia de \texttt{sub\_code}.
  \item \textbf{Derivadas temporales}: spreads y ratios entre features clave (\texttt{al$-$am}, \texttt{cg$-$by}, \texttt{s$-$t}); medias EWM con spans adaptados al horizonte (3,5 ó 5,7).
  \item \textbf{Normalización cross-seccional} por \texttt{ts\_index}: volatilidad, z-score, rango y distancia a mínimo/máximo.
  \item \textbf{Momentum}: diferencia y cambio porcentual; fuerza de tendencia e indicadores de volatilidad.
  \item \textbf{Estacionalidad}: módulos temporales 7/30/90; senos y cosenos de ciclo anual y de período 100; transformaciones $\log$ y cuadrado de \texttt{ts\_index}; interacciones horizonte$\times$tiempo.
  \item \textbf{Rolling y rezagos}: medias y desviaciones con ventanas 3/7/14; rezagos 1/3/7/14 de \texttt{feature\_al}; estadísticas históricas de \texttt{y\_target} por \texttt{code}, \texttt{sub\_code} y \texttt{sub\_category} (calculadas solo hasta $\texttt{ts\_index} \leq 3\,500$).
\end{enumerate}

La regularización se escala al SNR estimado por horizonte: \texttt{num\_leaves} entre 70--90, $\lambda_2$ entre 8--12, \texttt{min\_child\_samples} entre 150--250 y \texttt{early\_stopping} entre 160--200 rondas; con \texttt{feature\_fraction}=0.6 y \texttt{bagging\_fraction}=0.7. Las predicciones finales son el promedio de 5 semillas [42, 2024, 777, 1337, 9999].

% -----------------------------------------------------------
\subsection{Experimento 3: LightGBM con Polars, suavizado de outliers y selección de features}

El pipeline se articula en cuatro etapas:

\textbf{(1) Preprocesamiento.} El pipeline se reimplementa en \textbf{Polars} con reducción automática de tipos (Float64$\to$Float32, Int64$\to$Int8/16/32, String$\to$Categorical) para reducir uso de memoria. El \emph{forward-fill} causal combina train y test ordenados por grupo y tiempo, propagando solo features numéricas (nunca \texttt{y\_target} ni \texttt{weight}). Un \textbf{filtro Pseudo-Hampel} (mediana rolling $w=50$, umbral $5\sigma$) suaviza los picos extremos de 12 features ruidosas.

\textbf{(2) Ingeniería de features.} Sobre las originales se añaden: features de signo categórico e interacciones sign$\times$magnitud; combinaciones aritméticas (suma, resta, producto, cociente) entre pares de features clave; estadísticos rolling robustos (mediana, máximo, mínimo, z-score) con ventana 1\,000; acumulados EMA y MACD en spans [10, 30]; y máximos/mínimos por fila entre grupos de features.

\textbf{(3) Selección de features en dos pasadas.} Primero, un modelo de ranking con early stopping en 50 rondas estima la importancia de cada feature. Segundo, se entrena el modelo final solo con las $N_h$ features más importantes por horizonte ($N_1=197$, $N_3=177$, $N_{10}=197$, $N_{25}=217$).

\textbf{(4) Entrenamiento.} Ensemble de 5 semillas [42, 2026, 1232006, 3107, 2709], con los cuatro horizontes entrenados en paralelo con \texttt{joblib} (2 workers).

% -----------------------------------------------------------
\subsection{Experimento 4: Stacking LightGBM + XGBoost + Ridge con pesos por recencia}

Este experimento introduce dos innovaciones respecto a los anteriores.

\textbf{Innovación 1 — Ponderación por recencia.} Los pesos originales de cada observación se amplifican según su posición en el tiempo, dando hasta $40.5\times$ más importancia a las observaciones más recientes:
\begin{equation}
  \tilde{w}_i = w_i \cdot \bigl(0.5 + 40.0 \cdot \hat{t}_i\bigr), \qquad
  \hat{t}_i = \frac{t_i - t_{\min}}{t_{\max} - t_{\min} + \varepsilon}
  \label{eq:recency}
\end{equation}

\textbf{Innovación 2 — Stacking con meta-modelo Ridge.} Se entrenan LightGBM y XGBoost por separado (5 semillas cada uno), y un meta-modelo Ridge ($\alpha=1$, pesos no negativos, 5-Fold sin barajado) aprende qué peso asignar a cada modelo usando sus predicciones fuera de muestra:
\begin{equation}
  \hat{y}_{\text{stack}} = w_{\text{LGB}} \cdot \hat{y}_{\text{LGB}} + w_{\text{XGB}} \cdot \hat{y}_{\text{XGB}}, \qquad w_{\text{LGB}}, w_{\text{XGB}} \geq 0
  \label{eq:ridge_stack}
\end{equation}

\textbf{Features adicionales.} Se incorporan: wavelets EMA (diferencias entre spans 3, 10, 25 y 50); indicador de régimen de tendencia triple; oscilador estocástico en ventana 25; z-scores de reversión a la media; y desvíos cross-seccionales por \texttt{sub\_code}. Los hiperparámetros principales son: LightGBM con \texttt{lr}=0.005, \texttt{num\_leaves}=1024, $\lambda_2$=10; XGBoost con \texttt{lr}=0.015, \texttt{max\_leaves}=1024, \texttt{max\_depth}=10, early stopping 300.

% -----------------------------------------------------------
\subsection{Experimento 5: LightGBM minimalista con identificadores jerárquicos como features}

Se adopta un pipeline simplificado. Features: los 86 originales más \texttt{feature\_al}$-$\texttt{feature\_am} y las medias de grupo por (\texttt{code}, \texttt{sub\_code}, \texttt{sub\_category}, \texttt{horizon}) de \texttt{feature\_al} y \texttt{feature\_am}. Los identificadores jerárquicos se pasan como tipo \texttt{category} a LightGBM (\texttt{lr}=0.035, \texttt{n\_estimators}=1202, \texttt{num\_leaves}=64, \texttt{min\_child\_samples}=100, \texttt{feature\_fraction}=0.75, early stopping 100 rondas), permitiéndole aprender efectos fijos por entidad de forma implícita. El umbral de validación es $\texttt{VAL\_THRESHOLD} = 3\,500$.

% -----------------------------------------------------------
\subsection{Experimento 6: LightGBM minimalista con umbral de validación ajustado y entrenamiento determinista}

Mismo pipeline del Experimento 5 con tres ajustes: (i)~el umbral de validación se desplaza a $\texttt{VAL\_THRESHOLD} = 3\,503$; (ii)~la tasa de aprendizaje se reduce a 0.030 y la fracción de features por árbol a 0.70; y (iii)~se fija \texttt{seed=5} con \texttt{deterministic=True}, \texttt{force\_row\_wise=True} y \texttt{num\_threads=1} para garantizar reproducibilidad bit-a-bit entre ejecuciones.

% -----------------------------------------------------------
\subsection{Experimento 7: Ensemble dual LightGBM — modelos por horizonte y por sub-categoría con mezcla ponderada por potencia}

Se entrena una segunda dimensión de especialización: 4 modelos por horizonte (Pool A, vista temporal) y 5 modelos por sub-categoría (Pool B, vista estructural), con la misma ingeniería de features del Experimento 5 más \emph{target encoding} de \texttt{sub\_code} (media de \texttt{y\_target} calculada desde $\texttt{ts\_index} \leq 3\,500$) y rezagos \texttt{feature\_al\_lag1}, \texttt{feature\_am\_lag1}. Hiperparámetros: \texttt{lr}=0.030, \texttt{n\_estimators}=1200, \texttt{num\_leaves}=64, \texttt{bagging\_fraction}=0.75, \texttt{bagging\_freq}=5, \texttt{random\_state}=101. Las predicciones de ambos pools se fusionan mediante mezcla ponderada por potencia con $p=2.5$:
\begin{equation}
  \hat{y} = \frac{\mathrm{score}_h^{p}}{\mathrm{score}_h^{p} + \mathrm{score}_s^{p}} \cdot \hat{y}_h
          + \frac{\mathrm{score}_s^{p}}{\mathrm{score}_h^{p} + \mathrm{score}_s^{p}} \cdot \hat{y}_s, \qquad p = 2.5
  \label{eq:power_blend}
\end{equation}

% ============================================================
% 5. EXPERIMENTOS Y RESULTADOS
% ============================================================

\section{Experimentos y Resultados}

% -----------------------------------------------------------
\subsection{Experimento MLP: Red neuronal densa con embeddings de entidad}

La Tabla~\ref{tab:mlp_results} reporta las pérdidas MSE durante el entrenamiento sobre el 25\,\% de entidades submuestradas (corte de validación $\texttt{ts\_index} > 3\,060$).

\begin{table}[h!]
  \caption{Pérdidas MSE — Experimento MLP.}
  \label{tab:mlp_results}
  \centering
  \begin{tabular}{lcc}
    \toprule
    Fase & MSE & Época \\
    \midrule
    Mejor pérdida de validación & 758.1 & 1 \\
    Pérdida final de entrenamiento (reentrenamiento completo) & 252.4 & 50 \\
    \bottomrule
  \end{tabular}
\end{table}

Las 1\,447\,107 predicciones sobre test tienen: media $= -4.07$, desv.~estándar $= 66.4$, mín. $= -2\,773$, máx. $= 3\,993$. Score público: \textbf{0.0634}.

% -----------------------------------------------------------
\subsection{Experimento LSTM: Red recurrente con embeddings de entidad}

El modelo totalizó 412\,315 parámetros entrenables y se entrenó durante 100 épocas con función de pérdida MAE sobre el 25\,\% de entidades submuestradas (corte de validación $\texttt{ts\_index} > 3\,060$). Score público: \textbf{0.0891}.

% -----------------------------------------------------------
\subsection{Experimento 1: LightGBM multi-horizonte con ingeniería de features causales}

El conjunto de entrenamiento comprende 5\,337\,414 observaciones ($\texttt{ts\_index} \in [1, 3601]$) y el test 1\,447\,107 ($\texttt{ts\_index} \in [3602, 4376]$). La Tabla~\ref{tab:target_stats} muestra las estadísticas descriptivas del objetivo por horizonte.

\begin{table}[h!]
  \caption{Estadísticas descriptivas de \texttt{y\_target} por horizonte (entrenamiento).}
  \label{tab:target_stats}
  \centering
  \begin{tabular}{ccccc}
    \toprule
    Horizonte & Media & Desv. estándar & Asimetría & Curtosis \\
    \midrule
    1  & $-0.083$ & $11.700$ & $ 5.28$ & $528.5$ \\
    3  & $-0.252$ & $19.361$ & $ 2.01$ & $243.8$ \\
    10 & $-0.776$ & $33.842$ & $ 0.83$ & $176.8$ \\
    25 & $-1.682$ & $52.823$ & $ 0.85$ & $141.5$ \\
    \bottomrule
  \end{tabular}
\end{table}

La Tabla~\ref{tab:horizon_scores} reporta los scores de validación local por horizonte. Score público: \textbf{0.2675}.

\begin{table}[h!]
  \caption{Score de validación local por horizonte — Experimento 1 (holdout $\texttt{ts\_index} > 3500$).}
  \label{tab:horizon_scores}
  \centering
  \begin{tabular}{ccc}
    \toprule
    Horizonte & Filas validación & Score local \\
    \midrule
    1  & 43\,460 & 0.07547 \\
    3  & 43\,023 & 0.13232 \\
    10 & 40\,967 & 0.23661 \\
    25 & 37\,812 & 0.29283 \\
    \midrule
    \textbf{Global} & --- & \textbf{0.25083} \\
    \bottomrule
  \end{tabular}
\end{table}

% -----------------------------------------------------------
\subsection{Experimento 3: LightGBM con Polars, suavizado de outliers y selección de features}

La Tabla~\ref{tab:horizon_scores_exp3} reporta los scores locales por horizonte y el número de features seleccionadas. Score público: \textbf{0.2712}.

\begin{table}[h!]
  \caption{Score de validación local por horizonte — Experimento 3.}
  \label{tab:horizon_scores_exp3}
  \centering
  \begin{tabular}{ccc}
    \toprule
    Horizonte & Features seleccionadas ($N_h$) & Score local \\
    \midrule
    1  & 197 & 0.08037 \\
    3  & 177 & 0.11295 \\
    10 & 197 & 0.20699 \\
    25 & 217 & 0.28280 \\
    \bottomrule
  \end{tabular}
\end{table}

% -----------------------------------------------------------
\subsection{Experimento 4: Stacking LightGBM + XGBoost + Ridge con pesos por recencia}

La Tabla~\ref{tab:stacking_scores} reporta los scores individuales y los pesos Ridge por horizonte. Score público: \textbf{0.2824}.

\begin{table}[h!]
  \caption{Scores locales y pesos Ridge por horizonte — Experimento 4.}
  \label{tab:stacking_scores}
  \centering
  \begin{tabular}{cccccc}
    \toprule
    $h$ & LGB local & XGB local & Stacking local & $w_{\text{LGB}}$ & $w_{\text{XGB}}$ \\
    \midrule
    1  & 0.08947 & 0.01490 & 0.08109 & 78.6\% & 21.4\% \\
    3  & 0.13313 & 0.03721 & 0.13573 & 64.4\% & 35.6\% \\
    10 & 0.24442 & 0.03942 & 0.24379 & 94.3\% &  5.7\% \\
    25 & 0.29130 & 0.15028 & 0.27791 & 90.6\% &  9.4\% \\
    \midrule
    Global & --- & --- & \textbf{0.24446} & --- & --- \\
    \bottomrule
  \end{tabular}
\end{table}

% -----------------------------------------------------------
\subsection{Experimento 5: LightGBM minimalista con identificadores jerárquicos como features}

La Tabla~\ref{tab:exp5_conv} reporta la iteración de mejor desempeño y el RMSE de validación por horizonte (holdout $\texttt{ts\_index} > 3500$). Score público: \textbf{0.3212}.

\begin{table}[h!]
  \caption{Convergencia por horizonte — Experimento 5.}
  \label{tab:exp5_conv}
  \centering
  \begin{tabular}{ccc}
    \toprule
    Horizonte & Mejor iteración & RMSE validación \\
    \midrule
    1  & 133 & 0.001134 \\
    3  &  93 & 0.001818 \\
    10 & 107 & 0.002676 \\
    25 & 117 & 0.002908 \\
    \bottomrule
  \end{tabular}
\end{table}

% -----------------------------------------------------------
\subsection{Experimento 6: LightGBM minimalista con umbral de validación ajustado y entrenamiento determinista}

La Tabla~\ref{tab:exp6_conv} reporta la convergencia por horizonte (holdout $\texttt{ts\_index} > 3503$). Score público: \textbf{0.3377}.

\begin{table}[h!]
  \caption{Convergencia por horizonte — Experimento 6.}
  \label{tab:exp6_conv}
  \centering
  \begin{tabular}{ccc}
    \toprule
    Horizonte & Mejor iteración & RMSE validación \\
    \midrule
    1  &  53 & 0.001142 \\
    3  & 118 & 0.001825 \\
    10 & 541 & 0.002653 \\
    25 & 433 & 0.002843 \\
    \bottomrule
  \end{tabular}
\end{table}

% -----------------------------------------------------------
\subsection{Experimento 7: Ensemble dual LightGBM — modelos por horizonte y por sub-categoría con mezcla ponderada por potencia}

La Tabla~\ref{tab:exp7_scores} reporta los scores de validación local de cada modelo. El peso promedio resultante es 67.1\% para el Pool\,B (sub-categoría) y 32.9\% para el Pool\,A (horizonte). Las 1\,447\,107 predicciones finales tienen media $-0.869$, desv.~estándar $7.440$, mínimo $-158.05$ y máximo $86.75$. Score público: \textbf{0.3429}.

\begin{table}[h!]
  \caption{Scores de validación local por modelo — Experimento 7.}
  \label{tab:exp7_scores}
  \centering
  \begin{tabular}{llc}
    \toprule
    Pool & Modelo & Score local \\
    \midrule
    \multirow{5}{*}{Por horizonte}
      & $h=1$   & 0.0731 \\
      & $h=3$   & 0.1289 \\
      & $h=10$  & 0.2510 \\
      & $h=25$  & 0.3271 \\
      & \textit{Promedio} & \textit{0.1950} \\
    \midrule
    \multirow{6}{*}{Por sub-categoría}
      & \texttt{DPPUO5X2} & 0.2126 \\
      & \texttt{NQ58FVQM} & 0.2367 \\
      & \texttt{PHHHVYZI} & 0.2879 \\
      & \texttt{PZ9S1Z4V} & 0.3133 \\
      & \texttt{V8BKY1IV} & 0.2203 \\
      & \textit{Promedio} & \textit{0.2541} \\
    \bottomrule
  \end{tabular}
\end{table}

La Tabla~\ref{tab:resumen_final} resume los scores públicos de todos los experimentos con submission en el leaderboard.

\begin{table}[h!]
  \caption{Ranking de scores públicos en el leaderboard.}
  \label{tab:resumen_final}
  \centering
  \begin{tabular}{clc}
    \toprule
    Exp. & Descripción & Score público \\
    \midrule
    MLP  & Red densa con embeddings, 25\% entidades       & 0.0634 \\
    LSTM & Red recurrente con embeddings, 25\% entidades  & 0.0891 \\
    1    & LightGBM, 183 features, 5 semillas              & 0.2675 \\
    3    & LightGBM + Polars + Hampel + interacciones      & 0.2712 \\
    4    & LGB + XGB + Ridge stacking, recency weighting   & 0.2824 \\
    5    & LightGBM minimalista, IDs jerárquicos, thr=3500 & 0.3212 \\
    6    & LightGBM minimalista, IDs jerárquicos, thr=3503 & 0.3377 \\
    \textbf{7} & \textbf{Dual ensemble: horizonte + sub-cat, mezcla $p^{2.5}$} & \textbf{0.3429} \\
    \bottomrule
  \end{tabular}
\end{table}

% ============================================================
% 6. DISCUSIÓN
% ============================================================

\section{Discusión}

\paragraph{Redes neuronales vs.\ LightGBM.}
Los experimentos MLP (0.0634) y LSTM (0.0891) fallaron frente a cualquier variante de LightGBM. El MLP alcanzó su mejor pérdida de validación en la primera época (MSE = 758.1) y no volvió a mejorar; sus predicciones de test presentaron una desviación estándar de 66.4, muy por encima de la del objetivo real ($\approx 35$), confirmando que el modelo no aprendió la distribución subyacente. Tres factores explican el fracaso: el submuestreo al 25\% de entidades reduce la señal disponible, la curtosis extrema del objetivo ($>500$ en $h=1$) domina la función MSE con outliers, y la ausencia de ingeniería de features temporales causales priva a la red de las transformaciones que sí aprovecha LightGBM. El primer experimento LightGBM (Exp.~1, 0.2675) superó a las redes en $3\times$--$4\times$ usando el conjunto completo, confirmando el patrón observado en M5 \citep{makridakis2022m5}: en datos tabulares de bajo SNR, LightGBM explota cientos de features mejor que las arquitecturas secuenciales.

\paragraph{La complejidad no escala; la identidad sí.}
Los pipelines más elaborados de Exp.~3 (Polars, Hampel, selección en dos pasadas) y Exp.~4 (stacking LGB+XGB+Ridge con recencia) aportaron mejoras modestas sobre Exp.~1 ($+0.0037$ y $+0.0149$ respectivamente). El stacking llegó a perjudicar en $h=1$ (0.08109 vs.\ 0.08947 del LGB solo) y $h=25$ (0.27791 vs.\ 0.29130), con XGBoost contribuyendo apenas 5.7\% de peso en $h=10$. El salto cualitativo llegó con Exp.~5 ($+0.0388$, de 0.2824 a 0.3212) mediante un pipeline radicalmente más simple: pasar \texttt{code}, \texttt{sub\_code}, \texttt{sub\_category} y \texttt{horizon} como variables de tipo \texttt{category}, permitiendo a LightGBM aprender efectos fijos implícitos por entidad. Los modelos convergieron en apenas 93--133 iteraciones, evidenciando que la señal clave era la identidad de cada serie, no la ingeniería de features temporales.

\paragraph{Sensibilidad al esquema de validación y especialización por sub-categoría.}
El Exp.~6 desplazó el umbral de validación en solo 3 observaciones ($3\,500 \to 3\,503$) y obtuvo $+0.0165$ en el score público; la convergencia para $h=10$ pasó de 107 a 541 iteraciones y para $h=25$ de 117 a 433. Esta sensibilidad a un cambio del 0.08\% en el eje temporal refleja el riesgo de sobreajuste al leaderboard cuando el SNR es muy bajo. El Exp.~7 añadió cinco modelos por sub-categoría (Pool B, promedio 0.2541) frente a cuatro por horizonte (Pool A, promedio 0.1950), resultando en un peso promedio de 67.1\% para Pool B. La mejora incremental fue modesta ($+0.0052$), lo que indica que la señal de \texttt{sub\_category} ya estaba implícitamente codificada en los identificadores categóricos de Exp.~6.

\paragraph{Patrón de dificultad por horizonte.}
En todos los experimentos LightGBM, $h=1$ es el horizonte más difícil (score $\approx 0.07$--$0.09$) y $h=25$ el más fácil ($\approx 0.28$--$0.33$), patrón que se corresponde directamente con la curtosis del objetivo (528.5 en $h=1$ frente a 141.5 en $h=25$). Las series de corto plazo están dominadas por ruido de alta frecuencia, mientras las de largo plazo exhiben componentes estructurales más estables y recuperables.

% ============================================================
% REFERENCIAS
% ============================================================

\bibliographystyle{abbrvnat}
\bibliography{references}


\end{document}